
\begin{center}
  \large
  \textbf{MODE KUASINORMAL DALAM RUANG-WAKTU SCHWARZSCHILD ANTI-DE SITTER}
\end{center}
\addcontentsline{toc}{chapter}{ABSTRAK}
% Menyembunyikan nomor halaman
\thispagestyle{empty}

\begin{flushleft}
  \setlength{\tabcolsep}{0pt}
  \bfseries
  \begin{tabular}{ll@{\hspace{6pt}}l}
  Nama Mahasiswa &:& Ilham Rasyid Machfudi\\
  NRP &:& 5001221124\\
  Departemen &:& Fisika FSAD-ITS\\
  Dosen Pembimbing &:& Dr. rer. nat. Bintoro Anang Subagyo\\
  \end{tabular}
  \vspace{4ex}
\end{flushleft}
\textbf{Abstrak}

% Isi Abstrak
Mode kuasinormal merupakan osilasi teredam karakteristik yang muncul ketika
lubang hitam dikenai perturbasi kecil, dengan frekuensi kompleks
$\omega = \omega_R - i\omega_I$ yang bagian realnya menyatakan frekuensi osilasi
dan bagian imajinernya menyatakan laju peluruhan. Penelitian ini menentukan dan
membandingkan spektrum frekuensi mode kuasinormal perturbasi medan skalar pada
dua ruang-waktu, yaitu Schwarzschild yang datar secara asimtotik dan
Schwarzschild Anti-de Sitter. Frekuensi eigen diperoleh melalui metode Frobenius
yang dikombinasikan dengan metode pecahan berlanjut; khusus pada ruang-waktu
Schwarzschild Anti-de Sitter, relasi rekurensi multisuku yang dihasilkan
direduksi terlebih dahulu menjadi relasi tiga suku menggunakan eliminasi
Gaussian. Perhitungan dilakukan untuk medan skalar tak bermassa ($\mu = 0$) dan
bermassa ($\mu = 1$), momentum sudut $\ell = 0, 1, 2$, serta mode fundamental dan
overtone pertama. Hasil menunjukkan bahwa pada ruang-waktu Schwarzschild
peningkatan massa lubang hitam menurunkan bagian real maupun imajiner frekuensi,
sedangkan pada Schwarzschild Anti-de Sitter peningkatan jari-jari horizon
peristiwa $r_+$ justru menaikkan kedua komponen tersebut. Peningkatan momentum
sudut secara konsisten menaikkan bagian real dan menurunkan bagian imajiner
frekuensi pada kedua ruang-waktu, sedangkan peningkatan indeks overtone maupun
penambahan massa medan skalar menghasilkan kecenderungan yang berbeda antara
ruang-waktu Schwarzschild dan Schwarzschild Anti-de Sitter. Seluruh mode yang
diperoleh memiliki $\omega_I > 0$, sehingga lubang hitam Schwarzschild maupun
Schwarzschild Anti-de Sitter bersifat stabil terhadap perturbasi medan skalar,
baik untuk medan tak bermassa maupun bermassa.
\vspace{2ex}\\
\noindent \textbf{Kata Kunci: \emph{Medan Skalar, Mode Kuasinormal, SAdS, Schwarzschild}}